\documentclass[10pt,a4paper]{amsart}
\usepackage[margin=19mm]{geometry}
\usepackage{amsmath,amssymb,mathtools}
\usepackage[scr=rsfs,cal=euler]{mathalfa}
\usepackage{xfrac}
\usepackage[unicode,hidelinks,bookmarks=false]{hyperref}
\newcommand{\ie}{i.e.\ }
\newcommand{\cf}{cf.\ }
\newcommand{\resp}{respectively\ }
\newcommand{\Subsection}{\subsection}
\providecommand{\nobreakdash}{\nobreak\hbox{-}}
\setlength{\emergencystretch}{2em}
\multlinegap0pt
\usepackage{bm}
\usepackage{bbm}
\usepackage{dsfont}
\usepackage{xspace}
\usepackage{cancel}
\usepackage{mathtools}
\usepackage{amsthm}
\makeatletter
\@ifundefined{theorem}{\newtheorem{theorem}{Theorem}}{}
\makeatother
\makeatletter
\expandafter\ifx\csname Thm\endcsname\relax
\newenvironment{Thm}[1][]{\refstepcounter{alphabet}%
	\bigskip%
	\noindent%
	{\bf Theorem \Alph{alphabet}}%
	\ifthenelse{\equal{#1}{}}{}{ (#1)}%
	{\bf .} \itshape}{\vskip 8pt}
\fi
\expandafter\ifx\csname pf\endcsname\relax
\newenvironment{pf}[1][]{%
	\vskip 3mm
	\noindent
	\ifthenelse{\equal{#1}{}}%
	{{\slshape Proof. }}%
	{{\slshape #1.} }%
}%
{\qed\bigskip}
\fi
\expandafter\ifx\csname rems\endcsname\relax
\newenvironment{rems}{%
	\bigskip
	\noindent \textsl{{\sl Remarks. }}}{\bigskip}
\fi
\makeatother
\title[Strichartz estimates for higher order Schr\"odinger equation]{Strichartz estimates for higher order Schr\"odinger equations with Partial regular initial data}
\author{Vishvesh Kumar, Shyam Swarup Mondal, Iswarya Sitiraju, Manli Song}
\date{Source: arXiv:2508.15670v1 (CC BY 4.0)}
\begin{document}
\maketitle

\noindent\emph{Source and licence.} Strichartz estimates for higher order Schr\"odinger equations with Partial regular initial data. Version arXiv:2508.15670v1 (CC BY 4.0), distributed under the Creative Commons Attribution 4.0 International licence, as recorded in the arXiv licence metadata for this exact version and shown on the abstract page https://arxiv.org/abs/2508.15670v1. Full rights notice: licenses/arxiv-2508.15670-CC-BY-4.0.txt.\par\smallskip

\noindent\textit{Editorial scope.} Counted unit: the Theorem whose source label is \texttt{None} in the article (source0.tex lines 953--960, source label \texttt{None}), delivered together with the article's own complete original proof of it (source0.tex lines 961--1053, one \texttt{proof} environment of 93 lines). No mathematics is written, repaired or re-derived: statement and proof are the article's own text line for line. Required context retained uncounted: nonlin (lines 276--278); higher order (lines 307--312); eq11111 (lines 888--890); inhom (lines 896--898); evospace (lines 934--936); Banach1 (lines 944--945); Banach2 (lines 946--948). Every definition, display and label of those lines is retained unchanged. Omitted from this package and not redistributed: 1--275, 279--306, 313--887, 891--895, 899--933, 937--943, 949--952, 1054--1714. Nothing inside a retained excerpt is omitted or altered: every retained body is delivered exactly as the article's lines are written, so no omission receipt of any kind accompanies this package. Citations: the retained excerpts carry no citation command at all, so no bibliography entry of the article is transcribed and no reference is redistributed. Reference adaptation: none was needed and none was made; every reference inside the delivered unit resolves inside it, so no unresolved reference or citation is produced. Printed slips kept literal: no printed word, symbol, hypothesis or number of the counted statement or of its proof was altered. Layout and class: the article's own class is \texttt{amsart} with packages including amssymb, amsthm, amsmath, multirow, enumerate, ifthen, graphicx, tkz-euclide, hyperref, amsfonts, amscd, mathabx, color, amsbsy; the wrapper is the lane's pinned \texttt{amsart} editorial wrapper at 10pt on A4, which restates the theorem-like environments the retained text needs and adds the article's own macro definitions that the retained text needs (none), with extra wrapper packages (bm, bbm, dsfont, xspace, cancel, mathtools). The delivered numbering is the unit's own. Assets: no retained span contains a figure environment, a graphics-inclusion command, a PDF-page inclusion, a table or a TikZ picture; no image file is copied into this package, the package declares no asset dependency and no image file is redistributed with it. Licence: version arXiv:2508.15670v1 of this article is distributed under the Creative Commons Attribution 4.0 International licence, as recorded in the arXiv licence metadata for this exact version and shown on the abstract page https://arxiv.org/abs/2508.15670v1. A full rights notice is delivered at licenses/arxiv-2508.15670-CC-BY-4.0.txt. Characters: every retained excerpt is pure ASCII, so the delivered engine, XeLaTeX with its default Latin Modern text font, reports no missing character.
	\begin{equation} \label{nonlin}
	    \left|F_p(u)\right| \lesssim|u|^p \quad \text { and } \quad|u|\left|F_p^{\prime}(u)\right| \sim\left|F_p(u)\right| .
	\end{equation}

	\begin{align}\label{higher order}
		\begin{cases}
i\partial_t u (t, x) + (-\Delta)^{m/2} u (t, x) = F_p(u)(t, x), \\
u(0, x) = f(x),
\end{cases}
	\end{align}

			\begin{align}\label{eq11111}
				\left\|e^{i t (-\Delta)^{\frac{\alpha}{2}}} f\right\|_{L_t^q\left(\mathbb{R} ; L_x^r\left(\mathbb{R}^{d-k} ; L_y^{\tilde{r}}\left(\mathbb{R}^k\right)\right)\right)} \lesssim \|f\|_{\dot{H}_{x, y}^s}
			\end{align}

			\begin{align}\label{inhom}
				\left\|\int_0^t e^{i(t-\tau)  (-\Delta)^{\frac{\alpha}{2}}} F_p(u) d \tau\right\|_{L_t^q\left(\mathcal{I} ; L_x^r\left(\mathbb{R}^{d-k} ; L_y^{\tilde{r}}\left(\mathbb{R}^k\right)\right)\right)}  \leq C\left\|F_p\right\|_{L_t^{q_1'}\left(\mathcal{I} ; L_x^{r_1'}\left(\mathbb{R}^{d-k} ; L_y^{\tilde{r_1}'}\left(\mathbb{R}^k\right)\right)\right)}
			\end{align}

	\begin{equation} \label{evospace}
	    X(T)=\left\{u \in C\left([0, T] ; L_x^2 H_y^s\right) \cap L_t^q\left([0, T] ; L_x^r W_y^{s, \widetilde{r}}\right):  \|u\|_{L_t^q\left([0, T] ; L_x^r W_y^{s, \tilde{r}}\right)} \leq A\right\}
	\end{equation} with the distance

	\begin{align}\label{Banach1}
		\|N u\|_{X(T)} & \lesssim  \left\|f\right\|_{L_x^2\left(\mathbb{R}^{d-2} ; H_y^s\left(\mathbb{R}^2\right)\right)}+\|u\|_{X(T)}^p  \end{align}

	and \begin{align}\label{Banach2}
		d(u, v)=\|Nu-Nv\|_{L_t^q\left(I ; L_x^r L_y^{\tilde{r}}\right)}& \lesssim\|u-v\|_{X(T)}\left[  \|u\|_{X(T)}^{p-1}+\|v\|_{X(T)}^{p-1}\right],
	\end{align}

	\begin{theorem}
		For   $0<s \leq 1, m>2$ and $d \geq 3$,  let   $1<p<1+\frac{2m}{ d-2s}$ and $f \in$ $L_x^2\left(\mathbb{R}^{d-2} ; H_y^s\left(\mathbb{R}^2\right)\right)$.  Further, assume that  
		$(q, r, \tilde{r})$ with $2 \leq \widetilde{r} \leq r<\infty$ and $ 2<q \leq \infty $ satisfies $$			\frac{{{m}}}{q}=(d-2)\left(\frac{1}{2}-\frac{1}{r}\right)+2\left(\frac{1}{2}-\frac{1}{\widetilde{r}}\right).$$ Then there exists  $T>0$ 
		$$
		u \in C_t\left([0, T] ; L_x^2\left(\mathbb{R}^{d-2}; H_y^s(\mathbb{R}^{2})\right)\right) \cap   L_t^q\left([0, T] ; L_x^r \left(\mathbb{R}^{d-2}; W_y^{s, \widetilde{r}}(\mathbb{R}^{2})\right)\right)
		$$
		that satisfies  the higher order nonlinear schr\"odinger equation  (\ref{higher order}).
	\end{theorem}

	\begin{proof}
		
		 
		In order to prove the existence of a local solution to the nonlinear equation \eqref{higher order}, it is enough to show that the map $$N: X(T)\to Nu(t, x):=e^{i t (-\Delta)^{\frac{m}{2}}} f(x)-i \int_0^t e^{i(t-\tau) (-\Delta)^{\frac{m}{2}}} F_p(u)(\tau, x)\; d \tau$$
		is defines a contraction map on $X(T)$ defined in \eqref{evospace}. To prove $N$ is a contraction map on $X(T),$ by denoting $\mathcal{I}= [0,T]$,  our first step is to show  that
		\begin{align*}
			\left\|\left\langle\nabla_y\right\rangle^s F_p(u)\right\|_{L_t^{q_1^{\prime}}\left(\mathcal{I} ; L_x^{r_1^{\prime}} L_y^{\tilde{r}_1^{\prime}}\right)} \leq C T^{\beta(N, p, s)}\|u\|_{L_t^{q_1}\left(\mathcal{I} ; L_x^{r_1} W_y^{s, \tilde{r}_1}\right)}^p
		\end{align*}
		for $0<s \leq 1$ and $1<p<1+\frac{2m}{d-2s}$ and for  some  $\beta(N, p, s)>0$. 
		
		
		
		For $d \geq 3$, $0<s \leq 1$ and for a fixed exponent $p$ with  $1<p<1+\frac{2m}{d-2s}$, we   choose   $\varepsilon>0$ and  a pair $\left(q_1, r_1, \widetilde{r}_1\right)$ such that it satisfies the following conditions  
		$$
		\frac{(1-s)(p-1)}{2}<\varepsilon<\min \left\{\frac{d-2}{m^2}\left(1+\frac{2m}{d-2}-p\right), \frac{p-1}{2}\right\}
		$$
		and 
		\begin{equation} \label{mxm}
		    \frac{1}{r_1}=\frac{1}{p+1}, \quad \frac{1}{\widetilde{r}_1}=\frac{1}{2}-\frac{\epsilon}{p+1}, \quad \frac{1}{q_1}=\frac{d-2}{2m}-\frac{d-2}{m(p+1)}+\frac{m\varepsilon}{2(p+1)}.
		\end{equation}
		Note that above   $\left(q_1, r_1, \widetilde{r}_1\right)$ also satisfies the admissible condition 	$\frac{{{m}}}{q_1}+\frac{d-2}{r_1}+\frac{2}{\widetilde{r}_1}=\frac{d}{2}.$
        
        Using \eqref{nonlin}  and the H\"older inequality with \eqref{mxm}, we get the estimate 
		\begin{align}\label{Holder1}\nonumber
			\left\|F_p(u)\right\|_{L_t^{q_1^{\prime}}\left(\mathcal{I} ; L_x^{r_1^{\prime}} L_y^{\tilde{r}_1^{\prime}}\right)} & \lesssim\left\|\chi_{\mathcal{I}}(t)|u|^{p-1} | u|\right\|_{L_t^{q_1^{\prime}}\left(\mathcal{I} ; L_x^{r_1^{\prime}} L_y^{L_1^{\prime}}\right)} \\
			&\leq C T^{1-\frac{p+1}{q_1}}\|u\|_{L_t^{q_1}\left(\mathcal{I} ; L_x^{r_1} L_y^{(p-1)(p+1) / 2 \varepsilon}\right)}^{p-1}\|u\|_{L_t^{q_1}\left(\mathcal{I} ; L_x^{r_1} L_y^{\tilde{r_0}}\right)}.
		\end{align}
	In a similar manner, for $0<s \leq 1$, using the following fractional chain rule for Laplacian on $\mathbb{R}^d$ $$\||\nabla|^s F(u)\|_{L^c}\leq \||F'(u)\|_{L^a}\||\nabla|^sF(u)\|_{L^c},$$
    where $\frac{1}{c}=\frac{1}{a}+\frac{1}{b},$ along with the H\"older inequality,  we obtain
		\begin{align}\label{Holder2}
			\left\|\left|\nabla_y\right|^s F_p(u)\right\|_{L_t^{q_1^{\prime}}\left(\mathcal{I} ; L_x^{r_1^{\prime}} L_y^{\tilde{r}_1}\right)} \leq C T^{1-\frac{p+1}{q_1}}\|u\|_{L_t^{q_1}\left(\mathcal{I} ; L_x^{r_1} L_y^{(p-1)(p+1) / 2 \varepsilon}\right)}^{p-1}\|u\|_{L_t^{q_1}\left(\mathcal{I} ; L_x^{r_1} W_y^{s, \tilde{r}_1}\right)}.
		\end{align}
		Using the fact that  $W^{s, \widetilde{r}_1} \subseteq L^{\widetilde{r}_1}$ for $s \geq 0$ and  by  the Sobolev embedding in two dimensions
		$$
		W^{s, \widetilde{r}_1} \hookrightarrow W^{1-\frac{2 \varepsilon}{p-1}, \widetilde{r}_1} \hookrightarrow L^{\frac{(p-1)(p+1)}{2 \epsilon_1}},
		$$
		the inequalities \eqref{Holder1} and \eqref{Holder2}  further   reduces to 
		\begin{align}\label{using embedding }
			\left\||\nabla_y|^s F_p(u)\right\|_{L_t^{q_1^{\prime}}\left(\mathcal{I} ; L_x^{r_1^{\prime}} L_y^{\vec{r}^{\prime}}\right)} \leq C T^{\beta_1(d, p, s)}\|u\|_{L_t^{q_1}\left(\mathcal{I} ; L_x^{r_1} W_y^{s, \bar{r}_1}\right)}^p
		\end{align}
		with $1 \geq  s \geq 0$ and $\beta_1(d, p, s)=1-\frac{p+1}{q_1} =\frac{d-2}{2m}\left(1+\frac{2m}{d-2}-p\right)-\frac{m\varepsilon}{2}>0.$
		Thus using the inequalities \eqref{Holder1} and \eqref{using embedding }, we can write 
		\begin{align}\label{For banach}\nonumber
			\left\|\left\langle\nabla_y\right\rangle^s F_p(u)\right\|_{L_t^{q_1^{\prime}}\left(\mathcal{I} ; L_x^{r_1^{\prime}} L_y^{\tilde{r}_1^{\prime}}\right)} &\lesssim\left\|F_p(u)\right\|_{L_t^{q_1^{\prime}}\left(\mathcal{I} ; L_x^{r_1^{\prime}} L_y^{\tilde{r}_1}\right)}+\left\|\left|\nabla_y\right|^s F_p(u)\right\|_{L_t^{q_1^{\prime}}\left(\mathcal{I} ; L_x^{r_1^{\prime}} L_y^{\tilde{r}_1^{\prime}}\right)} \\
			&\lesssim  T^{\beta_1(d, p, s)}  \|u\|_{L_t^{q_1}\left(\mathcal{I} ; L_x^{r_1} W_y^{s, \bar{r}_1}\right)}^p.
		\end{align}
		In a similar way, we can have
		\begin{align}\label{To prove 2}\nonumber
			& \left\|F_p(u)-F_p(v)\right\|_{L_t^{q_1^{\prime}}\left(\mathcal{I}; L_x^{r_1^{\prime}} L_y^{\tilde{r_1}'}\right)} \\
			& \leq C T^{\beta_1(d, p, s)}\left(\|u\|_{L_t^{q_1}\left(\mathcal{I}; L_x^{r_1} W_y^{s, \tilde{r}_1}\right)}^{p-1}+\|v\|_{L_t^{q_1}\left(\mathcal{I} ; L_x^{r_1} W_y^{s, \tilde{r}_1}\right)}^{p-1}\right)\|u-v\|_{L_t^{q_1}\left(\mathcal{I}; L_x^{r_1} L_y^{\tilde{r}_1}\right)},
		\end{align} 
        by using the following elementary properties for the nonlinearity $F_p(u):$
        \begin{equation}
            |F_p(u)-F_p(v)| \leq C (|u|^{p-1}+|v|^{p-1}) |u-v|.
        \end{equation}
		Now our aim is to establish \eqref{Banach1} and \eqref{Banach2}. 
		By Plancherel's theorem, a dual version of  the Strichartz inequality \eqref{eq11111} with $s=0$, and from the inequality  \eqref{For banach}, one can see that 
		$$
		\begin{aligned}
			\sup _{t \in \mathcal{I}}\|Nu \|_{L_x^2 H_y^s} & \leq 
			\sup _{t \in \mathcal{I}} 	\|	e^{i t (-\Delta)^{\frac{m}{2}}} f \|_{L_x^2 H_y^s} +	\sup _{t \in \mathcal{I}}\left \| \int_0^t e^{i(t-\tau) (-\Delta)^{\frac{m}{2}}} F_p(u)\; d \tau \right\|_{L_x^2 H_y^s}\\ 		
			& \leq C\|f\|_{L_x^2 H_y^s}+C	\sup _{t \in \mathcal{I}} \left\|\left\langle\nabla_y\right\rangle^s \int_0^t e^{-i \tau (-\Delta)^{\frac{m}{2}}} F_p(u) d \tau\right\|_{L_x^2 L_y^2} \\
			& \leq C\|f\|_{L_x^2 H_y^s}+C\left\|\left\langle\nabla_y\right\rangle^s F_p(u)\right\|_{L_t^{q_1^{\prime}}\left(\mathcal{I} ; L_x^{r_1^{\prime}} L_y^{\tilde{r}_1^{\prime}}\right)}\\
			&	\leq C\|f\|_{L_x^2 H_y^s}+C	T^{\beta_1(d, p, s)}  \|u\|_{L_t^{q_1}\left(\mathcal{I} ; L_x^{r_1} W_y^{s, \bar{r}_1}\right)}^p,
		\end{aligned}
		$$
		for some $\left(q_1, r_1, \widetilde{r}_1\right) $  that satisfies 		$\frac{{{m}}}{q_1}+\frac{d-2}{r_1}+\frac{2}{\widetilde{r}_1}=\frac{d}{2}.$  Thus for $u \in X$, it follows that
		\begin{equation}\label{final1}
		    \sup _{t \in \mathcal{I}}\|Nu \|_{L_x^2 H_y^s} 	\leq C\|f\|_{L_x^2 H_y^s}+C	T^{\beta_1(d, p, s)} A^p.
		\end{equation}	 
		
		Let us now proceed to prove our next aim to show inequality of the form \eqref{Banach2}.  	 Now using  
		\eqref{inhom} and \eqref{To prove 2}, we get 
		\begin{align}\label{final2}\nonumber
			d(Nu, Nv) & =\|Nu-Nv\|_{L_t^q\left(I ; L_x^r L_y^{\tilde{r}}\right)}\\\nonumber
			&=\left\|\int_0^t e^{i(t-\tau)  (-\Delta)^{\frac{m}{2}}} \left( F_p(u)- F_p(v) \right)d \tau\right\|_{L_t^q\left(\mathcal{I} ; L_x^r L_y^{\tilde{r}}\right)} \\\nonumber
			& \leq C\left\|F_p(u)-F_p(v)\right\|_{L^{q_1^{\prime}}\left(\mathcal{I} ; L_x^{r_1^{\prime}} L_y^{\bar{r}_1^{\prime}}\right)}\\\nonumber
			&\leq C T^{\beta_1(d, p, s)}  \left(\|u\|_{L_t^{q_1}\left(\mathcal{I}; L_x^{r_1} W_y^{s, \tilde{r}_1}\right)}^{p-1}+\|v\|_{L_t^{q_1}\left(\mathcal{I} ; L_x^{r_1} W_y^{s, \tilde{r}_1}\right)}^{p-1}\right)\|u-v\|_{L_t^{q_1}\left(\mathcal{I}; L_x^{r_1} L_y^{\tilde{r}_1}\right)}\\ 				
			& \leq 2 C T^{\beta_1(d, p, s)} A^{p-1} d(u, v)
		\end{align} 	for $u, v \in X(T)$.  By choosing $A=2 C\|f\|_{L_x^2 H_y^s}$ and taking  sufficiently small $T>0$  such  that
		$$
		2 C T^{\beta_1(d, p, s)} A^{p-1} \leq \frac{1}{2}.
		$$
		Then \eqref{final1} and \eqref{final2}	 reduces to 	  
		\begin{align}\label{final11}
			\|Nu\|_{L^q\left(\mathcal{I} ; L_x^r W_y^{s, \bar{r}}\right)} \leq A
		\end{align}
		and  	 \begin{align}\label{final12}
			d(Nu, Nv)  \leq \frac{1}{2} d(u, v),
		\end{align}
		respectively 	for $u, v\in X(T).$  Thus from (\ref{final11}), we see that  $Nu \in X(T)$, i.e.,  $N$ maps $X(T)$ into itself. 
		Also from (\ref{final12}),  it follows that the map $N$ is a contraction mapping on the ball $\mathcal{B}(R)$ around the origin in the Banach space $X(T)$. Therefore, by Banach's fixed point theorem,  there exists a uniquely determined fixed point $u^*$ of the operator $N$, which means $u^*=Nu^* \in X(T)$.  This fixed point $u^*$ will be  our mild solution to (\ref{higher order})    on $[0, T]$. This implies that there exists a local   small data Sobolev solution $u^*$ of the equation $ u^*=Nu^* $ in $ X(T)$, which also gives the solution to the   equation (\ref{higher order})  and this completes the proof of the local well-posedness theorem. 
	\end{proof}

\end{document}
